Finger-specific motor intent is a clinically meaningful target for post-stroke neurorehabilitation, where residual muscle activity may be weak, variable, and patient-specific. This paper studies five-finger intent decoding from impaired-arm high-density surface electromyography (sEMG) using \physiomio, a bilateral longitudinal dataset collected from stroke patients~\citep{physiomioData2026}. We formulate the task as multilabel prediction and evaluate recurrent, convolutional, and graph-based decoders under a common signal-processing protocol: label alignment, 20--450 Hz Butterworth filtering, Symlet-4 wavelet denoising, 200 ms overlapping windows, and twelve handcrafted time- and frequency-domain descriptors per channel-window pair. In the direct baseline comparison, the LSTM obtains the highest subset accuracy, also called exact match ratio (0.545), while the GNN obtains the strongest macro F1 (0.706) and macro AUPRC (0.776). We then evaluate optimized CNN students and 1D ResNet teachers as deployment-oriented alternatives for Raspberry Pi 5-supported rehabilitation hardware~\citep{psrrepo2026,raspberrypi52023}. CNN-Large gives the strongest single-split aggregate CNN result, with 0.593 subset accuracy, 0.798 finger accuracy, 0.714 macro F1, 0.881 macro AUROC, and 0.812 macro AUPRC. For hardware deployment, CNN-Micro is selected because it retains similar aggregate performance at 158K parameters and an estimated 154 KB INT8 footprint. Healthy-to-impaired transfer improves the CNN-Base branch but gives mixed LSTM results. Overall, the study shows that compact neural decoders can convert post-stroke impaired-arm sEMG into five-finger intent estimates while making model-size and deployment constraints explicit.
